\section{Avoiding fixed modifications in the small-cut construction}
\label{sec:high-density-protected-cuts}

The small-cut construction can avoid any fixed set of vertices.
This permits added edges to be removed after the construction.

\begin{theorem}[Protected vertices]
\label{thm:high-density-protected-cyclic-sets}
Fix $k\ge2$, $m\ge0$, and $0<\varepsilon\le1/2$.
Put $C=\lceil C_\varepsilon\rceil$ and
\[
 B_{k,m,\varepsilon}
 =(k+m+1)(9C+8k+4m+9).
\]
Let $X$ be a normalized $k$-pole of rank $h\ge B_{k,m,\varepsilon}$.
Let $Z\subseteq V(X)$ satisfy $|Z|\le m$.
If some terminal pair has at least $(1/2+\varepsilon)2^h$ paths, then $X$ contains two disjoint vertex sets $V_1,V_2$ with these properties:
\begin{enumerate}
\item Each $V_i$ avoids $Z$ and every port of $X$, and contains a cycle.
\item In every cubic completion of $X$, each $V_i$ has boundary at most four.
\end{enumerate}
In particular, either set gives a cyclic cut of size at most four in every completion.
\end{theorem}

\begin{proof}
Use the normalization and series decomposition in Theorem~\ref{thm:high-density-pair-small-cyclic-cut}.
Call the number of bridgeless pieces $q$.
Each piece retains normalized terminal-path density at least $1/2+\varepsilon$.

If $q\ge k+m+2$, exclude the first and last pieces.
Also exclude pieces containing a restored port or a vertex of $Z$ that survives normalization.
At most $(k-2)+m$ further pieces are excluded.
At least two whole pieces remain.
Use their original vertex sets, without subdivision vertices on boundary edges.
They avoid $Z$ and all ports, contain cycles, and have boundary at most two in every completion.
A protected vertex that was suppressed is one of the restored ports and was already accounted for.

If $q\le k+m+1$, some piece has rank
\[
 b\ge\frac h{k+m+1}\ge9C+8k+4m+9.
\]
Its terminal distance is at least $b-C$.
There are at most $2C$ internal layers with more than two vertices.
Mark these layers and the at most $2(k-2)$ layers meeting endpoints of subdivided edges.
Also mark the layers of protected vertices that survive normalization.
This adds at most $m$ layers.

Partition the internal distance layers into disjoint groups of four.
The number of unmarked groups is at least
\[
 \left\lfloor\frac{b-C-1}{4}\right\rfloor-2C-2(k-2)-m
 \ge6.
\]
Choose two such groups and apply Lemma~\ref{lem:four-thin-layers-short-cycle}.
The resulting cycles have length at most four, avoid $Z$, and use no subdivided edge.
Internal layers avoid the two retained terminals, and unmarked layers avoid all restored ports.
Their disjoint vertex sets remain unchanged when the ports are restored.
Each has boundary at most four in every cubic completion.
\end{proof}

\begin{corollary}[Avoiding fixed added edges]
\label{cor:high-density-forbidden-edges}
Let $X$ be a normalized $k$-pole, and let $E_0$ be at most $s$ designated edges of $X$.
If $X$ has a terminal pair of density at least $1/2+\varepsilon$ and
\[
 \beta(X)\ge(k+2s+1)(9\lceil C_\varepsilon\rceil+8k+8s+9),
\]
then the two cyclic vertex sets in Theorem~\ref{thm:high-density-protected-cyclic-sets} avoid every endpoint of every edge in $E_0$.
Deleting the edges of $E_0$ and changing incidences at their endpoints leaves both cycles and both boundaries unchanged.
The last statement applies when all other vertex incidences remain unchanged.
\end{corollary}

\begin{proof}
Protect the at most $2s$ endpoints of the designated edges.
Neither a selected vertex nor an edge incident with it is affected by the stated changes.
\end{proof}

\begin{corollary}[Unconditional augmented inequalities in a five-pole]
\label{cor:five-pole-unconditional-added-edge}
Let $Y$ be a normalized five-pole of rank $h$ inside a cubic graph $G$ with no cyclic cut of size at most four.
For distinct ports $i,j,k,l$, write
\[
 a_{kl}=2^{-h}P_{kl}(Y),\qquad
 U_{ijkl}=2^{-h}F_{\{i,j,k,l\}}(Y),\qquad
 u_{ij|kl}=2^{-h}f_{\{i,j,k,l\},ij|kl}(Y).
\]
Then, uniformly over the choices of ports,
\[
 a_{kl}+U_{ijkl}-u_{ij|kl}\le1+o(1)
 \qquad(h\to\infty).
\]
No positive lower bound on $P_{ij}(Y)$ is required.
\end{corollary}

\begin{proof}
Add an edge $ij$ to $Y$, obtaining a normalized three-pole $T$ of rank $h+1$.
The exact added-edge identity gives
\[
 P_{kl}(T)
 =2^h\bigl(a_{kl}+U_{ijkl}-u_{ij|kl}\bigr).
\]
Suppose the parenthesized expression exceeds $1+\eta$ for some fixed $\eta>0$ along arbitrarily large ranks.
Then the normalized $k$--$l$ path density in $T$ exceeds $1/2+\eta/2$.
Apply Theorem~\ref{thm:high-density-protected-cyclic-sets} with protected vertices $i,j$.
It supplies two disjoint cyclic vertex sets contained in $Y-\{i,j\}$.

The selected sets avoid $i,j$ and the three remaining ports of $T$.
Every changed incidence is therefore outside both selected sets.
Thus their boundary sizes in $G$ equal those in any cubic completion of $T$.
Both boundaries have size at most four.
Either selected set therefore gives a forbidden cyclic cut in $G$.
This contradiction proves the assertion for every fixed $\eta>0$.
There are only finitely many port choices, so the error is uniform.
\end{proof}
